\begin{center}\begin{tikzpicture}[font=\small,>=stealth,line join=round]
\draw[gray!40,rounded corners=3pt] (0,0) rectangle (14.2,8.8);\node at (7.1,8.200000000000001) {底数范围有原因：a=1变成水平线，负底数不能全取实数};
\node at (3.1500000000000004,7.3) {$1^x=1$};\begin{scope}[shift={(2.8512195121951223,3.4981132075471697)},x=1.1951219512195124cm,y=0.6226415094339622cm]\draw[->,gray] (-1.8,0)--(2.3,0) node[right]{$x$};\draw[->,gray] (0,-0.8)--(0,4.5) node[above]{$y$};\draw[gray,thick,solid,domain=-1.5:2,samples=180] plot(\x,{1});\fill[orange!85!black] (0,1) circle(2pt);\fill[orange!85!black] (1,1) circle(2pt);\node[below left=5pt] at (0,0){$0$};\node[below=5pt] at (1,0){$1$};\end{scope}\node[align=center,text width=5.8cm] at (3.1500000000000004,1.7) {每处高度都为1\\单列为常数函数};\node at (9.95,7.3) {$(-2)^x$：只标整数示例};\begin{scope}[shift={(9.651219512195123,3.4981132075471697)},x=1.1951219512195124cm,y=0.6226415094339622cm]\draw[->,gray] (-1.8,0)--(2.3,0) node[right]{$x$};\draw[->,gray] (0,-0.8)--(0,4.5) node[above]{$y$};\fill[orange!85!black] (0,1) circle(2pt);\fill[orange!85!black] (2,4) circle(2pt);\node[below left=5pt] at (0,0){$0$};\node[below=5pt] at (0.5,0){$\frac12$};\node[below=5pt] at (2,0){$2$};\draw[gray,densely dotted] (.5,0)--(.5,3.5);\end{scope}\node[align=center,text width=5.8cm] at (9.95,1.7) {$x=\frac12$时无实数值\\不能连成全实数上的指数曲线};
\node[text width=12.6cm,align=center] at (7.1,.65) {两图同尺度；右图整数点只是示例，并不表示负底数幂的全部可用指数。};
\end{tikzpicture}\end{center}
